\subsection{Solar panel temperature sensors}\label{sec:temp_sensors}
There are four SFCS50-3-3-TB-5 temperature sensors \cite{temp_probe_datasheet} connected to the system to log temperatures of the solar panels \ref{eis:1.3}. The sensors are connected to an ADAM-4015 Data Aquisition Module \cite{adam_datasheet}, which is connected to an adapter that converts the Modbus RTU to USB so that it can be connected to the ODROID (See \Cref{fig:temperature_sensors_connection}).

\begin{figure}[!ht]
  \centering
    \includegraphics[width=0.6\textwidth]{images/ODROID_temperature_sensors.pdf}
      \caption{The connection of the weather station}\label{fig:temperature_sensors_connection}
\end{figure}

\subsubsection{Casting and scaling of temperature values}
Temperatures received from the ADAM-4015 module are encoded as a \verb|uint16_t|, but it was unknown how the module was programmed so it was yet to be determined what this value represents. To overcome this issue, The ``AdamApax .NET Utility'' \cite{AdamApax} was installed to investigate the settings that the module was operating at. The settings that were found can be seen in \Cref{fig:ADAM_-50-150}.\\

\begin{figure}[!ht]
  \centering
    \includegraphics[width=0.75\textwidth]{images/ADAM_-50-150.png}
      \caption{Settings of the ADAM-4015 Data Aquisition Module. Input range (top) shows the temperature range of -50\degree C to 150\degree C. The fourth connection shows the value that is used (23823), other sensors are not connected.}\label{fig:ADAM_-50-150}
\end{figure}

It was likely that the \verb|uint16_t| values could be cast to \verb|int|, and then be scaled between -50\degree C and 150\degree C. This would result in a temperature of $\frac{23823}{2^{16}-1} \cdot 200-50 = 22.70$\degree C. To confirm this, the settings were changed to output between -200\degree C and 200\degree C. The results can be found in \Cref{fig:ADAM_-200-200}.\\

\begin{figure}[!ht]
  \centering
    \includegraphics[width=0.75\textwidth]{images/ADAM_-200-200.png}
      \caption{Alternative settings of the ADAM-4015 Data Aquisition Module. Input range (top) shows the temperature range of -200\degree C to 200\degree C. The fourth connection shows the value that is used (36491), other sensors are not connected.}\label{fig:ADAM_-200-200}
\end{figure}

This capture was done close after the image in \Cref{fig:ADAM_-50-150} and results in a temperature of $\frac{36491}{2^{16}-1} \cdot 400-200 = 22.73$\degree C, which is similar enough to the previous calculation to conclude that this calculation is indeed correct.\\

The temperature values are therefore cast using the function found in Script \ref{castTempData}, after which they are scaled using the loop found in Script \ref{scaleTempData}.

\scriptsize
\begin{lstlisting}[language=C,caption={Function used to cast temperature data},label={castTempData}]
double* CastTempData(uint16_t *temperatures){

    //Cast the uint16_t array with temperatures to double
    double *double_temperatures;
    double_temperatures = (double *) calloc(NB_SENSORS, sizeof(double));

    int i;
    for(i = 0; i < NB_SENSORS; i++){
        double_temperatures[i] =(double)((int)((int16_t)temperatures[i]));
    }

    return double_temperatures;
}
\end{lstlisting}
\normalsize

\scriptsize
\begin{lstlisting}[language=C,caption={Function used to scale temperature data},label={scaleTempData}]
for(i = 0; i < NB_SENSORS; i++){
    double tmp = cast_temperatures[i]/65535*200-50; //scale linearly between -50 and 150 degrees
    scaled_temperatures[i] = roundf(tmp*100)/100; //round on two decimals
}
\end{lstlisting}
\normalsize

\subsubsection{Hotswap}
Due to the fact that the ADAM module gives values for ports that are available but not connected (65535) it is possible to implement a hotswap functionality to the temperature sensors. This means that it is possible to remove or add temperature sensors while the system is operational. The value of 65535 corresponds to 150\degree C and is not considered a realistic value for solar panel temperatures. This value will be sent to the server but will not be logged. It will be discarded instead. Whenever a sensor is added, however, this value will change and the server will know that a temperature sensor was added. Similarly, it will know when a sensor was disconnected by checking if the temperature value equals 150\degree C.

\subsubsection{Robustness}
A couple of checks are done in \verb|readTemperatureSensors| to ensure that the system can run 24/7 \ref{eis:2.3}. This is done by returning zero whenever a check fails (for example, a failed connection). The function that handles the data (see Script \ref{sec:code_overview}) will then check if a zero was returned. Whenever this happens, it will bypass updating the local display and the sending of data to the server.\\

The function will first check if a connection can be established to the ADAM module, this is done similar to the check in \verb|readWeatherStation| (see Script \ref{check_modbus_connection}). After reading the registers containing the temperature data, the function will check if all registers were read using the code found in Script \ref{check_temp_data_size}.

\scriptsize
\begin{lstlisting}[language=C,caption={Checking if the received data corresponds to the amount of requested registers},label={check_temp_data_size}]
if(rc!=NB_SENSORS){
	printf("FAILED (rc=%d) \n",rc);
	modbus_close(ctx);
	modbus_free(ctx);
	return 0;
}
\end{lstlisting}
\normalsize
